% Procedencia: COMPATIBILIDAD_HOLONOMIA_E_INFORMACION.md, secciones 2--9.
% Desarrollo científico conservado; gestión y correspondencia en archivo separado.
\section{Compatibilité constitutive, holonomie relative et information opérationnelle}
\label{md:comp:capitulo}

\subsection{Réciprocité différentielle de la vitesse et de l’impédance}
\label{md:comp:reciprocidad}

L’antécédent constitutif de la section~\ref{iii:sec:constitutiva} construit, sur les feuillets
conservés de sa représentation,
\begin{equation}
 \begin{gathered}
 T=e^{-xI-y\mathcal R},\qquad
 r_\pm=f\bigl(e^{-30(x\pm y)}\bigr),\\
 V=(I-T^{90})(I-T^{120})^{-1}=r_+P_++r_-P_-.
 \end{gathered}
 \label{md:comp:constitutiva}
\end{equation}
où
\[
 f(s)=\frac{1-s^3}{1-s^4},\qquad x>|y|.
\]
Les projecteurs de feuillet restent fixés. La chambre physique marquée
\(x>y>0\) et sa conjuguée sont contenues dans cette collection analytique.
Nous définissons les coordonnées logarithmiques normalisées
\begin{equation}
 c_\ell=\log\widehat c=-\log r_+-\log r_-,\qquad
 z_\ell=\log\widehat Z=\log r_--\log r_+.
 \label{md:comp:logaritmos}
\end{equation}
Ces lettres ne désignent ni de nouvelles constantes ni le contre-angle \(C^*\).
Soient
\[
 u=30(x+y),\qquad v=30(x-y),\qquad
 g(w)=\log f(e^{-w}),\qquad a=g'(u),\quad b=g'(v).
\]

\begin{proposition}[Réciprocité constitutive différentielle]
\label{md:comp:prop-reciprocidad}
Dans le domaine \(x>|y|\), on a
\begin{equation}
 g'(w)=\frac{3}{e^{3w}-1}-\frac{4}{e^{4w}-1}>0,
 \qquad
 D(c_\ell,z_\ell)=-30
 \begin{pmatrix}a+b&a-b\\a-b&a+b\end{pmatrix}.
 \label{md:comp:jacobiano}
\end{equation}
Les valeurs propres du jacobien sont \(-60a\) et \(-60b\). Il est donc
symétrique et défini négatif, son déterminant est \(3600ab>0\), et
\begin{equation}
 \boxed{\partial_y\log\widehat c=\partial_x\log\widehat Z.}
 \label{md:comp:reciprocidad-cruzada}
\end{equation}
\end{proposition}
\begin{proof}
On dérive les identités
\(c_\ell=-g(u)-g(v)\) et \(z_\ell=g(v)-g(u)\).
La positivité de \(g'\) provient de la décroissance stricte, sur \(k>0\),
de \(k/(e^{kw}-1)\). Elle résulte également de l’expression rationnelle positive
de la section~\ref{md:comp:restitucion}. La matrice obtenue possède les
valeurs propres et le déterminant indiqués ; l’égalité de ses entrées
non diagonales donne \eqref{md:comp:reciprocidad-cruzada}.
\end{proof}

L’égalité relie la réponse de la vitesse à la séparation
orientée et la réponse de l’impédance à la coordonnée commune.
Les deux sensibilités proviennent d’une même fonction spectrale. C’est une
réciprocité de cette collection constitutive ; elle n’est pas identifiée sans preuve
à une loi thermodynamique ou d’électromagnétisme macroscopique.

\subsection{Potentiel spectral et connexion de profondeur deux avec Catalan}
\label{md:comp:potencial}

\subsubsection{Construction du potentiel}
\label{md:comp:construccion-potencial}

À partir des canaux déjà construits, nous définissons
\begin{equation}
 \mathscr F(x,y)=\sum_{\sigma=\pm1}\sum_{n\geq1}
 \left[
 \frac{e^{-120n(x+\sigma y)}}{120n^2}
 -\frac{e^{-90n(x+\sigma y)}}{90n^2}
 \right].
 \label{md:comp:serie-potencial}
\end{equation}

\begin{proposition}[Potentiel des deux réponses logarithmiques]
\label{md:comp:prop-potencial}
La série \eqref{md:comp:serie-potencial} définit sur \(x>|y|\) un potentiel
qui satisfait
\begin{equation}
 \boxed{\partial_x\mathscr F=c_\ell,\qquad
        \partial_y\mathscr F=z_\ell.}
 \label{md:comp:gradiente-potencial}
\end{equation}
Sa hessienne est le jacobien de
\eqref{md:comp:jacobiano} ; par conséquent, \(\mathscr F\) est
strictement concave. La constante additive est fixée par
\(\mathscr F\longrightarrow0\) lorsque \(x-|y|\longrightarrow\infty\).
\end{proposition}
\begin{proof}
Sur tout compact du domaine \(x>|y|\), les quantités \(x\pm y\)
ont un minimum positif. Une puissance de \(n\) multipliée par une
exponentielle décroissante est sommable ; ainsi, la série et ses dérivées
de tout ordre fixé convergent uniformément sur ce compact.
En dérivant terme à terme et en utilisant
\(-\log(1-z)=\sum_{n\geq1}z^n/n\), on obtient les deux identités
de \eqref{md:comp:gradiente-potencial}. La hessienne coïncide alors
avec \eqref{md:comp:jacobiano}, qui est défini négatif.
La normalisation à l’infini est celle de la série présentée.
\end{proof}

\(\mathscr F\) est un potentiel mathématique de la réponse. On ne lui
attribue pas d’unités d’énergie et on ne le déclare pas potentiel thermodynamique.

\subsubsection{Réalisation par un même moment spectral}
\label{md:comp:momento-espectral}

Soit
\[
 D_2(z)=\sum_{n\geq1}\frac{z^n}{n^2}.
\]
Sa reconnaissance comme dilogarithme est postérieure à cette définition par
série. Pour \(|z|<1\), il satisfait
\((z\partial_z)^2D_2(z)=z/(1-z)\). Dans cette notation,
\begin{equation}
 \mathscr F=
 \sum_{\sigma\in\{+,-\}}
 \left[\frac{D_2(q_\sigma^{120})}{120}
       -\frac{D_2(q_\sigma^{90})}{90}\right].
 \label{md:comp:potencial-dos-momentos}
\end{equation}
Le moment orienté de Catalan construit dans la section~\ref{iii:sec:catalan} satisfait
\begin{equation}
 \mathcal C_2(s)=\operatorname{Im}D_2(is)
 =\sum_{k\geq0}\frac{(-1)^ks^{2k+1}}{(2k+1)^2},\qquad
 \mathcal C_2(1)=G_{\rm Cat}.
 \label{md:comp:catalan-momento}
\end{equation}

\begin{proof}
Les puissances paires de \(i\) sont réelles et les impaires ont une partie
imaginaire \((-1)^k\). Cela donne la série de
\eqref{md:comp:catalan-momento}, qui converge absolument même
en \(s=1\). Pour \(0\leq s<1\), appliquer terme à terme deux fois
\(s\partial_s\) redonne \(s/(1+s^2)\), le coefficient orienté de la
résolvante du quart de tour déjà présente dans le corpus.
L’identité dérivée en \(s=1\), si elle est utilisée, s’interprète
par prolongement analytique ou limite d’Abel : la série deux fois
dérivée ne converge pas au sens ordinaire à cette extrémité.
L’identification \(\mathcal C_2(1)=G_{\rm Cat}\) provient directement
de la série originale absolument convergente.
\end{proof}

Par conséquent, le moment de Catalan et le potentiel constitutif sont réunis
par un même moment spectral de profondeur deux, évalué sur des
supports différents : le quart de tour orienté et les contractions
des secteurs \(90/120\). Cette composition n’identifie pas Catalan
à \(\mathscr F\), ni n’affirme que sa seule valeur extrême détermine
tous les canaux. La collection analytique et ses arguments conservent
une information qu’une valeur isolée ne contient pas.

\subsubsection{Compatibilité avec l’amplification}
\label{md:comp:amplificacion-potencial}

Pour \(T_m=T\otimes I_{9^m}\) et
\(\tau_m=\operatorname{Tr}/(2\cdot9^m)\), on a exactement
\begin{equation}
 \mathscr F_m=2\tau_m\left[
 \frac{D_2(T_m^{120})}{120}
 -\frac{D_2(T_m^{90})}{90}\right]=\mathscr F.
 \label{md:comp:potencial-amplificado}
\end{equation}
La multiplicité \(9^m\) annule la normalisation. Le même argument
conserve les dérivées. Cela atteste l’amplification de la
représentation, non tout raffinement imaginable du TPK.
L’opérateur temporel n’est pas remplacé par une liste de valeurs propres
sans marques de feuillet.

Pour tout lacet lisse fermé dans cette collection analytique,
\begin{equation}
 \oint(c_\ell\,dx+z_\ell\,dy)=0.
 \label{md:comp:integral-cerrada}
\end{equation}
Cette égalité se rapporte à la projection constitutive à deux
coordonnées. Elle n’implique pas que l’histoire enrichie ou son holonomie
opératoire soient triviales.

\subsection{Restitution exacte et sensibilité finie}
\label{md:comp:restitucion}

\begin{proposition}[Restitution analytique des coordonnées angulaires]
\label{md:comp:prop-restitucion}
Avec \(L=\log(16/9)\), l’application
\((x,y)\longmapsto(c_\ell,z_\ell)\) est un difféomorphisme réel analytique
du domaine \(x>|y|\) sur
\begin{equation}
 0<c_\ell+z_\ell<L,\qquad 0<c_\ell-z_\ell<L.
 \label{md:comp:imagen-logaritmica}
\end{equation}
\end{proposition}
\begin{proof}
La fonction \(g\) croît strictement de \(\log(3/4)\) à \(0\).
Les combinaisons
\[
 c_\ell+z_\ell=-2g(u),\qquad c_\ell-z_\ell=-2g(v)
\]
récupèrent \(u,v\) au moyen de \(g^{-1}\). Alors
\(x=(u+v)/60\) et \(y=(u-v)/60\).
L’injectivité est maintenue lorsqu’on restreint à l’image effectivement
produite par HMT ; on ne déclare pas que HMT produise tout le cône.
\end{proof}

Avec \(s=e^{-w}\), la dérivée a la forme rationnelle
\begin{equation}
 a(w)=g'(w)=
 \frac{s^3(s^2+2s+3)}{(1+s+s^2)(1+s+s^2+s^3)},
 \qquad 0<a(w)<\frac12.
 \label{md:comp:derivada-racional}
\end{equation}
La fonction est strictement décroissante, avec \(a(0+)=1/2\), et
\(a(w)\sim3e^{-3w}\) à l’infini. Pour prouver la décroissance,
on utilise
\begin{equation}
 a'(w)=-\frac{9}{4\sinh^2(3w/2)}
        +\frac{16}{4\sinh^2(2w)}<0,
 \label{md:comp:decrecimiento}
\end{equation}
car \(k/\sinh(kw/2)\) décroît strictement sur \(k>0\).
En outre,
\begin{equation}
 \frac12e^{-3w}<a(w)<3e^{-3w}.
 \label{md:comp:cotas-exponenciales}
\end{equation}
La borne inférieure résulte de la multiplication des dénominateurs et de l’utilisation de
\begin{equation}
 \begin{split}
 &2(s^2+2s+3)-(1+s+s^2)(1+s+s^2+s^3)\\
 &\hspace{1cm}=(1-s)(5+7s+6s^2+3s^3+s^4)>0.
 \end{split}
 \label{md:comp:polinomio-cota}
\end{equation}
La borne supérieure résulte de
\[
 3(1+s+s^2)(1+s+s^2+s^3)-(s^2+2s+3)>0.
\]

Si \(m=30(x-|y|)\) et \(M=30(x+|y|)\), la norme de l’inverse du
jacobien est \(1/[60a(M)]\). Pour ne pas confondre l’application vectorielle avec
le potentiel scalaire, nous écrivons
\(\mathbf F=(c_\ell,z_\ell)\).

\begin{proposition}[Bornes finies de restitution]
\label{md:comp:prop-sensibilidad}
Sur une région convexe où les deux arguments \(u,v\) appartiennent
à \([m_0,M_0]\subset(0,\infty)\), on vérifie
\begin{equation}
 60a(M_0)\|p-q\|
 \leq\|\mathbf F(p)-\mathbf F(q)\|
 \leq60a(m_0)\|p-q\|.
 \label{md:comp:cotas-finitas}
\end{equation}
\end{proposition}
\begin{proof}
On intègre la hessienne de \(\mathscr F\), c’est-à-dire le jacobien
de \(\mathbf F\), sur le segment. Pour la borne inférieure, on utilise
\[
 -\langle\mathbf F(p)-\mathbf F(q),p-q\rangle
 \geq60a(M_0)\|p-q\|^2
\]
et l’inégalité de Cauchy--Schwarz. La borne supérieure utilise la
norme maximale du jacobien. Ce sont des bornes d’erreurs finies dans le
domaine spécifié.
\end{proof}

En \(y=0\) et lorsque \(x\longrightarrow\infty\), les deux directions
se contractent également : le conditionnement matriciel vaut un,
mais la norme de l’inverse croît comme \(e^{90x}/180\).
L’inversibilité exacte et la récupération robuste face à l’erreur absolue
sont donc des propriétés différentes. Lorsque \(|y|\longrightarrow x\)
avec \(x>0\) fixé, il n’y a pas de singularité différentielle : une dérivée tend
vers \(1/2\) et l’autre vers \(a(60x)>0\).

\subsection{Le quart de tour provient du cycle dodécaphasique}
\label{md:comp:cuarto-giro}

On retrouve la construction de \eqref{iii:eq:quarter-intertwiner}. Dans \(\mathbb R^{12}\),
soient
\[
 C_{12}e_j=e_{j+1\bmod12},\qquad
 \Pi_4=P_4^{\rm cyc}\frac{I-C_{12}^6}{2},
\]
et
\begin{equation}
 \begin{gathered}
 u_0=(1,0,-1,0,1,0,-1,0,1,0,-1,0)^{\mathsf T},\qquad
 v_0=C_{12}u_0,\\
 W=(u_0,v_0)/\sqrt6.
 \end{gathered}
 \label{md:comp:isometria-ciclica}
\end{equation}
Le symbole \(\Pi_4\) distingue le projecteur des coordonnées
d’action \(Q,P\). On vérifie
\begin{equation}
 \begin{gathered}
 W^{\mathsf T}W=I,\qquad WW^{\mathsf T}=\Pi_4,\qquad
 C_{12}W=WJ_0,\\
 J_0=\begin{pmatrix}0&-1\\1&0\end{pmatrix},\qquad J_0^2=-I.
 \end{gathered}
 \label{md:comp:entrelazador-ciclico}
\end{equation}
En particulier, \(C_{12}^3W=-WJ_0\). Le plan possède ainsi un opérateur d’entrelacement
explicite avec le cycle ; aucune rotation externe n’a été introduite pour
ajuster une réponse.

L’antécédent radial de \eqref{rec:eq:entrelazamiento} réalise
\[
 S_\zeta=\operatorname{diag}(e^{\zeta/2},e^{-\zeta/2}),
 \qquad J_\zeta=S_\zeta J_0S_\zeta^{-1}.
\]
Sur ce plan, nous définissons
\begin{equation}
 U_\zeta(\theta)=S_\zeta e^{-\theta J_0}S_\zeta^{-1}
 =\begin{pmatrix}
 \cos\theta&e^\zeta\sin\theta\\
 -e^{-\zeta}\sin\theta&\cos\theta
 \end{pmatrix}.
 \label{md:comp:transporte-plano}
\end{equation}
Le signe de \(\theta\) fixe l’orientation choisie. Le rapport engendré
\(\xi=C^*/A\), avec \(|\xi|<1\) et les deux angles exprimés dans la même
unité, satisfait \(\zeta=\operatorname{artanh}\xi\).
Dans l’orientation conjuguée,
\((\varepsilon,\xi)\mapsto(-\varepsilon,-\xi)\) conserve la section
positive d’action au moyen de l’opérateur de récupération orienté du corpus ;
\(\zeta\) change de signe.

\subsection{Lacet relatif de deux transports}
\label{md:comp:lazo}

Soient
\begin{equation}
 U=U_{\zeta_s}(\theta),\qquad
 V=U_{\zeta_r}(\vartheta),\qquad
 H_{\rm lazo}=UVU^{-1}V^{-1}.
 \label{md:comp:def-lazo}
\end{equation}
Les produits se lisent avec une action de droite à gauche.
La lettre \(\vartheta\) désigne ici la seconde phase, pour la distinguer
de la coordonnée régionale \(\varphi\). Posons
\[
 \delta=\zeta_s-\zeta_r,\qquad
 \chi=\sinh\delta\,\sin\theta\,\sin\vartheta.
\]

\begin{proposition}[Trace et spectre du lacet relatif]
\label{md:comp:prop-lazo}
La composition \eqref{md:comp:def-lazo} satisfait
\begin{equation}
 \boxed{\begin{gathered}
 \det H_{\rm lazo}=1,\qquad
 \operatorname{tr}H_{\rm lazo}=2+4\chi^2,\\
 H_{\rm lazo}=I\ \Longleftrightarrow\ \chi=0.
 \end{gathered}}
 \label{md:comp:traza-lazo}
\end{equation}
Ses valeurs propres sont
\begin{equation}
 \lambda_\pm=(\sqrt{1+\chi^2}\pm|\chi|)^2
 =\exp\bigl(\pm2\operatorname{arsinh}|\chi|\bigr).
 \label{md:comp:espectro-lazo}
\end{equation}
\end{proposition}
\begin{proof}
En conjuguant simultanément par \(S_{\zeta_r}^{-1}\), l’opérateur \(V\)
est \(R(-\vartheta)\), et \(U\) possède les termes
\(e^{\pm\delta}\sin\theta\). La multiplication directe donne, dans cette
référence,
\[
 UV-VU=-2\chi\operatorname{diag}(1,-1).
\]
Comme
\[
 H_{\rm lazo}-I=(UV-VU)U^{-1}V^{-1},
\]
on obtient l’équivalence avec \(\chi=0\). La multiplication, ou
l’identité de traces pour les matrices unimodulaires, donne
\(\operatorname{tr}H_{\rm lazo}=2+4\chi^2\).
Le polynôme caractéristique est
\(\lambda^2-(2+4\chi^2)\lambda+1\), et ses racines sont celles de
\eqref{md:comp:espectro-lazo}.
\end{proof}

C’est une holonomie relative de la composition déclarée.
Elle n’est identifiée ni à \(\Gamma_9\) ni à une mise à jour complète
du TPK par similitude de noms.

\subsubsection{Spécialisation entièrement déterminée par la paire conjuguée}
\label{md:comp:lazo-conjugado}

Soient \(U_+=U_\zeta(\pi/2)\) et
\(U_-=U_{-\zeta}(\pi/2)\). Les deux satisfont
\(U_+^2=U_-^2=-I\), mais
\begin{equation}
 \begin{aligned}
 U_+U_-&=-\operatorname{diag}(e^{2\zeta},e^{-2\zeta}),\\
 H_{\rm lazo}=(U_+U_-)^2
 &=\operatorname{diag}(e^{4\zeta},e^{-4\zeta}).
 \end{aligned}
 \label{md:comp:lazo-diagonal}
\end{equation}
En écrivant \(\xi=\tanh\zeta=C^*/A\), avec les deux angles exprimés
dans la même unité, on obtient
\begin{equation}
 \begin{aligned}
 H_{\rm lazo}
 &=\operatorname{diag}\left(
 \left[\frac{1+\xi}{1-\xi}\right]^2,
 \left[\frac{1-\xi}{1+\xi}\right]^2\right),\\
 \operatorname{tr}H_{\rm lazo}-2
 &=\frac{16\xi^2}{(1-\xi^2)^2}.
 \end{aligned}
 \label{md:comp:lazo-razon-angular}
\end{equation}
La formule de \(H_{\rm lazo}\) compose le quart de tour et les formes
conjuguées de la même paire. Son expression est exacte dans le rapport angulaire
engendré \(\xi\).

Sur la branche positive d’action du corpus,
\(A_{\deg}=1000\alpha\) et
\(C^*_{\deg}=2(\eta_{\rm ret}+\alpha)\), avec
\(\hbar_{\rm ret}=s_0\eta_{\rm ret}\). Par conséquent,
\begin{equation}
 \begin{aligned}
 \xi&=\frac{\eta_{\rm ret}+\alpha}{500\alpha},\\
 \rho_{\rm lazo}=e^{4\zeta}
 &=\left(\frac{501\alpha+\eta_{\rm ret}}
                 {499\alpha-\eta_{\rm ret}}\right)^2,
 \qquad 0<\eta_{\rm ret}<499\alpha.
 \end{aligned}
 \label{md:comp:multiplicador-lazo}
\end{equation}
Les coefficients \(501\) et \(499\) sont \(500\pm1\) ;
\(500\) provient du facteur \(1000/2\) de la paire angulaire et n’est pas un
nouveau paramètre. Les axes orientés étant conservés, la restitution est
\begin{equation}
 \boxed{
 \eta_{\rm ret}=\alpha\left[
 500\frac{\sqrt{\rho_{\rm lazo}}-1}
          {\sqrt{\rho_{\rm lazo}}+1}-1\right].}
 \label{md:comp:restitucion-accion}
\end{equation}
La dérivée
\[
 \frac{d\eta_{\rm ret}}{d\rho_{\rm lazo}}
 =\frac{500\alpha}
 {\sqrt{\rho_{\rm lazo}}(\sqrt{\rho_{\rm lazo}}+1)^2}
\]
est positive. La chambre \(\eta_{\rm ret}>0\) correspond à
\(\rho_{\rm lazo}>(501/499)^2\) ; le bord
\(\eta_{\rm ret}\longrightarrow499\alpha\) correspond à
\(\rho_{\rm lazo}\longrightarrow\infty\).
C’est une lecture inverse du coefficient adimensionnel d’action à partir
du transport relatif et de la coordonnée \(\alpha\) déjà construite.
Elle conserve l’échelle \(s_0\) comme antécédent : une transformation
unimodulaire de forme ne détermine pas à elle seule cette échelle dimensionnelle.
L’orientation conjuguée inverse \(\rho_{\rm lazo}\) et transporte
l’étiquette d’action ; elle ne transforme pas \(\eta_{\rm ret}\) en action négative.

\subsubsection{Itération du lacet et accumulation orientée}
\label{md:comp:iteracion}

Soit
\[
 \rho_{\rm lazo}=e^{4\zeta}
 =\left(\frac{1+\xi}{1-\xi}\right)^2.
\]
Pour tout entier \(n\),
\begin{equation}
 \begin{gathered}
 H_{\rm lazo}^n
 =\operatorname{diag}(\rho_{\rm lazo}^{n},\rho_{\rm lazo}^{-n}),
 \\
 Q_n=\rho_{\rm lazo}^{n}Q_0,\quad
 P_n=\rho_{\rm lazo}^{-n}P_0.
 \end{gathered}
 \label{md:comp:iteracion-diagonal}
\end{equation}
La preuve est la multiplication de matrices diagonales, étendue à
\(n<0\) au moyen de l’inverse. Le produit
\(Q_nP_n=Q_0P_0\) et l’aire symplectique sont conservés, tandis que l’état change
pour \(\zeta\ne0\) et \(X_0\ne0\). Dans le programme de quatre
opérations, les incréments de phase déclarés ont une somme nulle à chaque
répétition ; cette comptabilité des phases ne détermine pas le transport
total de l’état.

Pour \(\zeta\ne0\) et \(Q_0\ne0\), on récupère exactement
\[
 n=\frac{\log|Q_n/Q_0|}{4\zeta}.
\]
Si \(Q_0=0\), la condition \(P_0\ne0\) permet d’utiliser
\[
 n=-\frac{\log|P_n/P_0|}{4\zeta}.
\]
La lecture nécessite de connaître la préparation et de conserver une référence.
L’extension multiplicative
\(H_{\rm lazo}^{n+m}=H_{\rm lazo}^{n}H_{\rm lazo}^{m}\)
correspond à un incrément logarithmique additif \(4(n+m)\zeta\).
Inverser le lacet change le signe de cet incrément.

Ainsi, un programme fixe et répété admet une réponse qui enregistre
son itération. C’est une réalisation mathématique d’accumulation
opérationnelle dans la composition étudiée. On n’identifie pas \(n\)
à toute la retenue ou la mémoire de \(\Gamma_9\), et l’on ne déduit pas une
entropie nulle pour la dynamique complète à partir de la périodicité
de son programme de phases. Si l’amplitude initiale est inconnue et
si l’on ne dispose que de l’état final, \(n\) n’est pas déterminé par
cette lecture isolée.

\subsubsection{Relèvement explicite aux douze phases}
\label{md:comp:levantamiento}

Soient
\[
 \widetilde S_\zeta=I-\Pi_4+WS_\zeta W^{\mathsf T},\qquad
 L_\pm=\widetilde S_{\pm\zeta}C_{12}^3
       \widetilde S_{\pm\zeta}^{-1}.
\]
Sur \(\operatorname{Ran}\Pi_4\) agissent \(U_+,U_-\), respectivement ;
sur le complément, tous deux agissent comme \(C_{12}^3\). Par conséquent,
\begin{equation}
 L_+L_-L_+^{-1}L_-^{-1}
 =I-\Pi_4+WH_{\rm lazo}W^{\mathsf T}.
 \label{md:comp:levantamiento-lazo}
\end{equation}
L’action complémentaire s’annule. La trace de l’opérateur à douze
phases est \(10+\operatorname{tr}H_{\rm lazo}\).
L’entrelacement et la composition sont exacts et sont vérifiés
rationnellement à l’aide des vecteurs entiers \(u_0,v_0\).
Ce relèvement réunit l’antécédent dodécaphasique et le transport
elliptique ; il n’affirme pas que sa séquence soit à elle seule la mise à jour
autonome originale du TPK.

\subsubsection{Raffinement et limite}
\label{md:comp:limite}

L’amplification
\(H_{{\rm lazo},m}=H_{\rm lazo}\otimes I_{9^m}\) conserve
\[
 \frac{1}{9^m}\operatorname{Tr}H_{{\rm lazo},m}
 =\operatorname{Tr}H_{\rm lazo},
\]
ainsi que la norme et le spectre, hormis les multiplicités.
Avec les inclusions \(v\mapsto v\otimes e_0\), on a
\[
 H_{{\rm lazo},m+1}\iota_m=\iota_m H_{{\rm lazo},m}.
\]
L’opérateur uniforme définit dans la limite inductive hilbertienne une
extension bornée, de même norme. C’est la limite de cette
amplification ; il ne remplace pas un autre système de raffinements
du corpus.

\subsection{Ce qui distingue un changement de coordinatisation d’un effet relatif}
\label{md:comp:carta-efecto}

Le transport de coordonnées
\[
 T_n=S_{\zeta_{n+1}}R(-\theta_n)S_{\zeta_n}^{-1}
\]
se télescope :
\begin{equation}
 T_{N-1}\cdots T_0
 =S_{\zeta_N}R\left(-\sum_{n=0}^{N-1}\theta_n\right)
 S_{\zeta_0}^{-1}.
 \label{md:comp:telescopaje}
\end{equation}
Si la coordinatisation et la phase reviennent, on obtient \(I\). Le commutateur de la
section~\ref{md:comp:lazo} compose des évolutions de formes distinctes
dans une référence commune ; ce n’est pas ce changement passif.

Sous un changement commun de représentation \(P\), les opérateurs
\(U,V,H_{\rm lazo}\) sont conjugués par \(P\).
La trace et le spectre de \(H_{\rm lazo}\) restent invariants.
Si, en outre, \(X\mapsto PX\) et
\(g\mapsto P^{-\mathsf T}gP^{-1}\), on conserve
\(X^{\mathsf T}gX\) et toute lecture énergétique correctement
transportée. Ainsi, un \(H_{\rm lazo}\ne I\) ne disparaît pas par
renommage des coordonnées.

Cela constitue un protocole mathématique relationnel fermé.
Sa mise en œuvre expérimentale exige que les évolutions et leurs
inverses soient des opérations physiques disponibles et que le système,
la référence et le contrôleur possèdent une réalisation commune.
On n’affirme ici ni force nouvelle ni source d’énergie.

\subsubsection{Conservation de l’aire et conservation d’une même énergie}
\label{md:comp:area-energia}

\begin{proposition}[Absence d’une métrique positive commune]
\label{md:comp:prop-metrica}
Chaque quart de tour conserve sa propre métrique :
\(g_\zeta=\operatorname{diag}(e^{-\zeta},e^\zeta)\) pour \(U_+\),
et \(g_{-\zeta}\) pour \(U_-\). Lorsque \(\zeta\ne0\), il n’existe pas
de métrique positive commune.
\end{proposition}
\begin{proof}
Les conditions
\[
 G=\begin{pmatrix}a&b\\b&d\end{pmatrix}>0,
 \qquad U_+^{\mathsf T}GU_+=G
\]
imposent \(b=0\) et \(d=ae^{2\zeta}\).
La condition \(U_-^{\mathsf T}GU_-=G\) exige
\(d=ae^{-2\zeta}\). Toutes deux ne sont compatibles qu’avec \(\zeta=0\).
\end{proof}

Par conséquent, conserver l’aire symplectique n’équivaut pas à conserver
une même forme énergétique pendant toute composition de ces
transports. Dans la référence
\(g_r=S_r^{-\mathsf T}S_r^{-1}\), avec \(\omega_r>0\),
\(X\ne0\) et \(E_r(X)=\omega_r X^{\mathsf T}g_rX/2\), on a
\begin{equation}
 \begin{gathered}
 \frac{E_r(H_{\rm lazo}X)}{E_r(X)}
 =\frac{x^{\mathsf T}\overline{H}_{\rm lazo}^{\mathsf T}
                   \overline{H}_{\rm lazo}x}{x^{\mathsf T}x},\\
 \overline{H}_{\rm lazo}=S_r^{-1}H_{\rm lazo}S_r,\qquad
 x=S_r^{-1}X.
 \end{gathered}
 \label{md:comp:cociente-energia}
\end{equation}
Les extrêmes sont les carrés des valeurs singulières.
Dans le cas conjugué diagonal, ils sont \(e^{\pm8|\zeta|}\), tandis que
les dilatations de l’état sont \(e^{\pm4\zeta}\).
Cette différence d’exposants est nécessaire. Changer physiquement
d’opération exige de prendre en compte le contrôleur ; le gain d’une
lecture n’atteste pas une création d’énergie.

\subsubsection{L’ordre nécessite une lecture orientée}
\label{md:comp:orden-orientado}

Échanger \(U_+\) et \(U_-\) inverse \(H_{\rm lazo}\).
La trace et l’ensemble des gains extrêmes sont égaux pour
\(H_{\rm lazo}\) et \(H_{\rm lazo}^{-1}\). Ces scalaires détectent
l’effet du lacet, mais non son orientation. Deux entrées
indépendantes avec leurs réponses vectorielles complètes reconstruisent
\(H_{\rm lazo}\) ; une mesure scalaire ne suffit pas en général.

Dans le cas conjugué, avec des axes de référence marqués, soient
\[
 G_Q=\frac{E_r(H_{\rm lazo}e_Q)}{E_r(e_Q)},\qquad
 G_P=\frac{E_r(H_{\rm lazo}e_P)}{E_r(e_P)}.
\]
La lecture
\begin{equation}
 \mathcal O=\frac14\log(G_Q/G_P)=4\zeta
 \label{md:comp:lector-orientado}
\end{equation}
récupère l’orientation et change de signe lorsqu’on inverse le lacet.
Les axes et la métrique doivent être transportés ensemble sous un changement
de représentation. On distingue ainsi l’orientation sans l’attribuer
à une trace qui l’élimine.

\subsection{Information d’histoire et statistique d’occupation}
\label{md:comp:informacion}

Les produits
\[
 U_+U_-U_+^{-1}U_-^{-1}
 \qquad\text{et}\qquad
 U_+U_+^{-1}U_-U_-^{-1}
\]
contiennent les mêmes quatre opérations avec les mêmes multiplicités.
Le second vaut \(I\) ; le premier est \(H_{\rm lazo}\ne I\) pour
\(\zeta\ne0\). Par conséquent, toute statistique qui dépend
uniquement des multiplicités attribue le même résultat à deux
protocoles ayant des actions différentes sur l’état.

La conclusion exacte est que cette statistique ne constitue pas une
description suffisante pour prédire la réponse de la composition.
Shannon n’est pas réfuté : une distribution sur des séquences ordonnées
peut représenter l’ordre. On distingue l’information qui a été
conservée dans l’observable choisie.

\begin{definition}[Information opérationnelle d’une histoire]
\label{md:comp:def-informacion}
Dans cette représentation, nous proposons d’appeler information opérationnelle d’une
histoire la classe de son transport ordonné discernable par les
lectures admises. Soit \(\mathcal D\) une collection de détecteurs
définis sur des opérateurs. Elle peut inclure
\[
 D_{X,L}(H_{\rm lazo})=L(H_{\rm lazo}X),
 \qquad D_{\rm tr}(H_{\rm lazo})=\operatorname{tr}H_{\rm lazo},
\]
avec une préparation \(X\) et un lecteur linéaire \(L\).
Deux histoires sont équivalentes si
\[
 D(H_{{\rm lazo},1})=D(H_{{\rm lazo},2})
 \qquad\text{pour tout }D\in\mathcal D.
\]
\end{definition}

Élargir la collection peut séparer des classes auparavant indiscernables.
Avec des réponses vectorielles complètes sur une base, on reconstruit tout
opérateur linéaire ; avec la seule trace, il peut subsister
\(H_{\rm lazo}\sim H_{\rm lazo}^{-1}\).
On distingue ainsi le détecteur d’une réponse de l’état du détecteur
spectral de l’opérateur.

Cette définition n’identifie pas le transport à toute la généalogie
TPK : des histoires distinctes peuvent produire le même opérateur dans une
représentation. L’information latente dans d’autres feuillets, incidences et
prolongements est conservée dans leurs propres lecteurs.

La dimension statistique et la dimension opérationnelle sont compatibles.
Le taux de croissance combinatoire, ou entropie topologique, compte
la croissance asymptotique des configurations admissibles ; un taux
d’entropie de Shannon nécessite en outre une distribution.
Une application d’observation détermine quelles transformations peuvent
être reconstruites. Aucune des preuves précédentes ne déclare une entropie
thermodynamique universelle nulle. L’apport consiste à préciser quelle description
est suffisante pour une tâche de prédiction ou de restitution.

\subsection{Conséquences et nouvelles définitions proposées}
\label{md:comp:consecuencias}

Les résultats précédents réunissent cinq notions, avec leurs domaines
et leurs lecteurs conservés :
\begin{enumerate}
 \item \textbf{Potentiel spectral constitutif.}
 \(\mathscr F\), défini par les moments \(90/120\), a pour
 ses deux dérivées le logarithme de la vitesse normalisée et le
 logarithme de l’impédance normalisée. C’est une définition mathématique
 postérieure aux canaux, sans interprétation énergétique ajoutée.

 \item \textbf{Réciprocité constitutive différentielle.}
 L’égalité des sensibilités croisées démontrée dans la
 section~\ref{md:comp:reciprocidad} fournit une contrainte
 conjointe vérifiable, outre les valeurs.

 \item \textbf{Défaut de fermeture relatif.}
 \(\operatorname{tr}H_{\rm lazo}-2\) détecte le lacet non trivial.
 Il s’accompagne d’une lecture orientée, comme \(\mathcal O\), lorsqu’il
 faut distinguer l’ordre. La trace seule ne définit pas toute la mémoire.

 \item \textbf{Restitution opérationnelle.}
 C’est la récupération de l’action de l’histoire à partir d’une collection
 spécifiée d’observables. Elle se distingue de la restitution de deux
 canaux et de la reconstruction intégrale de l’état.

 \item \textbf{Marge de restitution.}
 La borne \(60a(M_0)\) de la section~\ref{md:comp:restitucion}
 quantifie la stabilité absolue sur un domaine, tandis que
 l’injectivité affirme seulement l’unicité exacte.
\end{enumerate}

Les coordonnées de masse, de fréquence, de rayons et de température du
développement précédent restent intactes. Ce développement ajoute
deux niveaux : la compatibilité différentielle entre lecteurs scalaires
et la mémoire d’ordre dans les compositions opératoires.
On n’identifie pas le facteur de dilatation de \(H_{\rm lazo}\) à
une masse nouvelle sans passer par le lecteur de masses correspondant.
